\documentclass[10pt,a4paper]{amsart}
\usepackage[margin=19mm]{geometry}
\usepackage{amsmath,amssymb,mathtools}
\usepackage[scr=rsfs,cal=euler]{mathalfa}
\usepackage{xfrac}
\usepackage[unicode,hidelinks,bookmarks=false]{hyperref}
\newcommand{\ie}{i.e.\ }
\newcommand{\cf}{cf.\ }
\newcommand{\resp}{respectively\ }
\newcommand{\Subsection}{\subsection}
\providecommand{\nobreakdash}{\nobreak\hbox{-}}
\setlength{\emergencystretch}{2em}
\multlinegap0pt
\usepackage{dsfont}
\usepackage{amssymb}
\usepackage{csquotes}
\newtheorem{theorem}{Theorem}
\newtheorem{prop}{Proposition}
\newtheorem{lemma}{Lemma}
\newcommand{\Ind}{\mathbb{1}}
\newcommand{\Mod}{\,\mathrm{mod}\,}
\newcommand{\R}{\mathds{R}}
\newcommand{\Z}{\mathds{Z}}
\newcommand{\e}{\mathrm{e}}
\newcommand{\eps}{\varepsilon}
\newcommand{\starsum}{\ \sideset{}{^*}\sum}
\title{A Bombieri-Vinogradov theorem for exponential sums over products of $k$ primes}
\author{P\'ea Bazin}
\date{Source: arXiv:2607.15137v2 (CC BY 4.0)}
\begin{document}
\maketitle

\noindent\emph{Source and licence.} A Bombieri-Vinogradov theorem for exponential sums over products of $k$ primes. Version arXiv:2607.15137v2 (CC BY 4.0), distributed by arXiv under the Creative Commons Attribution 4.0 International licence, http://creativecommons.org/licenses/by/4.0/, as recorded in the arXiv licence metadata for this exact version and shown on the abstract page https://arxiv.org/abs/2607.15137v2. Full licence notice: \texttt{licenses/arxiv-2607.15137-CC-BY-4.0.txt}.\par\smallskip

\noindent\textit{Editorial scope.} Counted unit: the article's theorem main (source0.tex lines 103--108). The article's complete original proof of it is delivered with it (source0.tex lines 380--408, one \texttt{proof} environment of 29 lines, which the article places away from the statement). No mathematics is written, repaired or re-derived: the statement and the proof are the article's own lines, in the article's own notation, and no retained line is substituted or re-flowed.  Retained uncounted as the source-local context the proof needs: lines 174--177; lines 229--234; lines 284--286.  Nothing was omitted from any retained span, so the package carries no omission record.  No retained line contains a citation, so the wrapper carries no bibliography and no bibliography entry.  No retained line contains a figure environment, an image-inclusion command or drawing code, so no image is delivered and the package declares no asset dependency.  Editorial formatting only, in addition: an AMS \texttt{amsart} preamble that restates the article's own declarations and macros used by the retained lines, namely the environments \texttt{theorem}, \texttt{prop}, \texttt{lemma} on separate counters, together with the standard packages those lines need.  The retained environments print in this document as Theorem 1, Proposition 1, Lemma 1 in order of appearance, whereas the article prints the same items with the numbering of its own full document.  The complete native source of the article as distributed in the arXiv e-print for this exact version is bundled unchanged as \texttt{sources/205361/source0.tex}.
\begin{prop}\label{xi-pik}
    For $2\le Q \le x, \theta\in\R$ and $k \le \log x,$ we have
    $$\Xi(f_k, x, Q,\theta) \ll \left(x^{1/2}Q^2 + x^{5/6}Q + x + xQ^2\theta^{1/2}\right)(\log x)^{O(1)},$$ where the implicit constants are absolute.
\end{prop}

\begin{lemma}\label{decomp}
    Let $a,q$ with $(a, q) = 1.$ Define for $\chi$ a Dirichlet character, $$\chi_{/t} : n \mapsto \Ind_{t|n}\chi(n/t),$$ so that $\left\{\chi_{/t} : t|q, \chi\Mod \frac{q}{t}\right\}$ forms a basis of $q-$periodic functions over $\Z.$ Then
    \begin{equation}\label{decomp-eq}
        F(x; \e_{a/q}) = \frac{1}{\phi(q)}\sum_{st = q} \mu(s) t F(x; t, 0) + O\left(\frac{1}{\phi(q)}\sum_{\substack{q'st = q \\ q'\ne 1}} t\sqrt{q'}\starsum_{\chi\Mod q'} \left|F(x; \chi_{/t})\right| \right).
    \end{equation}
\end{lemma}

\begin{lemma}\label{sw-pik}
    There is an absolute constant $c$ such that for $A, \eps > 0,$ $k \le (2-\eps)\log\log x, $ $\lambda \le (\log x)^{-10}$ and $\chi$ a Dirichlet character of conductor $q \ll (\log x)^A,$ we have $$\Pi_k(x; \chi\e_\lambda) \ll_{\eps, A} x\left(q^{1/2}|\lambda|^{1/4-\eps} + \e^{-c\sqrt{\log x}}\right).$$ 
\end{lemma}

\begin{theorem}\label{main}
    Let $\eps, C > 0.$ There exists $B$ depending only on $C$ such that when $Q \le x^{1/3}(\log x)^{-B}, \theta \le Q^{-3}(\log x)^{-B}$  and $1 \le k \le (2-\eps)\log\log x,$ we have
    \begin{equation}\label{main-eq}
        \sum_{q\le Q} \max_{(a,q) = 1} \max_{|\lambda|\le\theta} \left|\sum_{n\le x} f_k(n) \e\left(\left(\frac{a}{q}+\lambda\right)n\right) - \frac{1}{\phi(q)}\sum_{st = q} \mu(s) t\sum_{\substack{n\le x \\ t|n}} f_k(n)\e(\lambda n) \right| \ll \frac{x}{(\log x)^C}.
    \end{equation}
\end{theorem}

\begin{proof}[Proof of Theorem \ref{main}]
    We apply Lemma \ref{decomp} to $f = f_k\e_\lambda$ to bound the left-hand side of (\ref{main-eq}) is bounded by
    \begin{align}
        \ll E :&= \sum_{q\le Q} \frac{1}{\phi(q)} \sum_{\substack{q'st = q \\ q'\ne 1}} t\sqrt{q'}\starsum_{\chi\Mod q'} \max_{|\lambda|\le \theta} \left|\Pi_k(x; \chi_{/t}\e_{\lambda})\right| \notag
        \\&= \sum_{t\le Q} t \sum_{1 < q'\le Q/t} \sqrt{q'} \starsum_{\chi\Mod q'}  \sum_{s\le Q/q't} \frac{1}{\phi(q'st)} \max_{|\lambda|\le \theta} \left|\Pi_k(x; \chi_{/t}\e_{\lambda})\right| \notag
        \\&\le \log Q \sum_{t\le Q} \frac{t}{\phi(t)} \sum_{1 < q'\le Q/t} \frac{\sqrt{q'}}{\phi(q')} \starsum_{\chi\Mod q'}  \max_{|\lambda|\le \theta} \left|\Pi_{k-\Omega(t)}(x/t; \chi\e_{\lambda t})\right|, \label{step99}
    \end{align}
    noting that we always have $\phi(q'st) \ge \phi(q')\phi(s)\phi(t)$ and $$\Pi_k(x; \chi_{/t}\e_\lambda) = \sum_{\substack{n \le x \\t|n \\ \Omega(n) = k}} \chi(n/t) \e(\lambda n) = \sum_{\substack{m\le x/t \\ \Omega(m) = k - \Omega(t)}} \chi(m) \e(\lambda tm) = \Pi_{k-\Omega(t)}(x/t;\chi\e_{\lambda t}).$$

    We now split the sum over $q'$ in (\ref{step99}) into dyadic intervals. For all $1\le t\le Q$ and $Q' \le Q/t,$ we have by Proposition \ref{xi-pik}
    \begin{align}
        \sum_{Q' < q'\le 2Q'} \frac{\sqrt{q'}}{\phi(q')} \starsum_{\chi\Mod q'} \max_{|\lambda|\le \theta} \left|\Pi_{k-\Omega(t)}(x/t; \chi\e_{\lambda t})\right| &\le Q'^{-1/2} \log Q' \big(\Xi(f_{k-\Omega(t)}, x/t,2Q', \theta t) - \Xi(f_{k-\Omega(t)}, x/t,Q', \theta t)\big) \notag\\&\ll \left(\frac{x^{1/2}Q'^{3/2}}{t^{1/2}} + \frac{x^{5/6}Q'^{1/2}}{t^{5/6}} + \frac{xQ'^{-1/2}}{t} + \frac{xQ'^{3/2}\theta^{1/2}}{t^{1/2}}\right) (\log x)^A \notag 
        \\&\ll \left(\frac{x^{1/2}Q^{3/2}}{t^{2}} + \frac{x^{5/6}Q^{1/2}}{t^{4/3}} + \frac{xQ'^{-1/2}}{t} + \frac{xQ^{3/2}\theta^{1/2}}{t^{2}}\right) (\log x)^A \notag
        \\&\ll \frac{x}{t} (\log x)^A \left(Q'^{-1/2} + (\log x)^{-B/2}\right)
        \label{step100}
    \end{align} 
    where $A$ is an absolute constant and we used the conditions $Q \le x^{1/3}(\log x)^{-B},$ $\theta \le Q^{-3}(\log x)^{-B}.$

    Thus, for $Q' \ge (\log x)^{B/7},$ the left-hand side of (\ref{step100}) is $\ll \frac{x}{t}(\log x)^{A-B/14}.$ 
    Moreover, when $Q'\le (\log x)^{B/7},$ we also have by Lemma \ref{sw-pik} with $\theta \le (Q't)^{-3}(\log x)^{-B}$
    \begin{align*}
        \sum_{Q' < q'\le 2Q'} \frac{\sqrt{q'}}{\phi(q')} \starsum_{\chi\Mod q'} \max_{|\lambda|\le \theta} \left|\Pi_{k-\Omega(t)}(x/t; \chi\e_{\lambda t})\right| &\ll Q'^{3/2} \frac{x}{t}\left(Q'^{1/2}(\theta t)^{1/4-\eps} + \e^{-c\sqrt{\log x}}\right)
        \\&\ll  \frac{x}{t}\left((\log x)^{-B/14 + \eps} + (\log x)^{B/4}\e^{-c\sqrt{\log x}}\right)
        \\&\ll \frac{x}{t}(\log x)^{A-B/14}.
    \end{align*}
    %shows the left-hand side of (\ref{step100}) is also $\ll Q'^{3/2}\frac{x}{t}\e^{-c\sqrt{\log x}} \ll \frac{x}{t}(\log x)^{A-B/2}$ when $Q'\le (\log x)^B.$ %\theta t \le Q^{-2}  
    Thus, we can inject this bound in (\ref{step99}) to get 
    $$E \ll (\log Q)^2 \sum_{t\le Q} \frac{t}{\phi(t)} \frac{x}{t} (\log x)^{A-B/2} \ll x(\log x)^{3+A-B/14}$$ as desired fixing $B = 14(A+C+3).$
\end{proof}

\end{document}
